\begin{figure}[tb]
\centering
\begin{tikzpicture}[x=1cm,y=1cm,>=Stealth,font=\small]
  \draw[gray!65,->] (-.4,0)--(10.7,0) node[right] {time};
  \draw[gray!65,->] (0,-1.55)--(0,1.7) node[above] {position};
  \fill[linkblue!4] (0,-1.2)--(10,-.2)--(10,2.2)--(0,1.2)--cycle;
  \fill[linkblue!17] (0,-.4)--(10,.6)--(10,1.4)--(0,.4)--cycle;
  \draw[linkblue,thick] (0,0)--(10,1);
  \draw[linkblue!50,dashed] (0,-1.2)--(10,-.2);
  \draw[linkblue!50,dashed] (0,1.2)--(10,2.2);
  \draw[linkblue!80] (0,-.4)--(10,.6);
  \draw[linkblue!80] (0,.4)--(10,1.4);
  \draw[densely dotted] (5,-1.15)--(5,1.8);
  \fill (5,.66) circle (2pt)
    node[right=4pt,above=1pt] {\(b\)};
  \fill (5,.30) circle (2pt)
    node[left=4pt,below=1pt] {\(b'\)};
  \node[anchor=west,fill=white,inner sep=2pt] at (7.5,1.75)
    {parent high};
  \node[anchor=west,fill=white,inner sep=2pt] at (7.8,1.02)
    {parent low};
  \node[anchor=north] at (5,-1.28) {one child time bin};
  \node[anchor=south west] at (.6,.15) {parent centerline};

  \node at (5,-2.10) {Position cross-section at the marked bin};
  \begin{scope}[shift={(5,-3.55)},x=1.3cm]
    \filldraw[fill=linkblue!4,draw=linkblue!50,dashed]
      (-3,-.42) rectangle (3,.42);
    \filldraw[fill=linkblue!10,draw=linkblue!65]
      (-1.6,-.30) rectangle (2.4,.30);
    \filldraw[fill=linkblue!17,draw=linkblue!80]
      (-1,-.18) rectangle (1,.18);
    \draw[gray!65,densely dotted] (0,-.90)--(0,.50);
    \draw[linkblue!65,densely dotted] (.4,-.30)--(.4,.30);
    \fill (.4,0) circle (2pt);
    \fill (-.5,0) circle (2pt);
    \node[anchor=north] at (.4,-.48) {\(b\)};
    \node[anchor=north] at (-.5,-.48) {\(b'\)};
    \node[anchor=north,font=\scriptsize,text=gray!80] at (0,-.96)
      {parent center};
    \node at (-2.25,.80) {parent high};
    \draw[linkblue!50] (-2.25,.59)--(-2.55,.42);
    \node at (1.10,.80) {intermediate about \(b\)};
    \draw[linkblue!65] (1.10,.59)--(1.65,.30);
    \node at (-2.20,-.72) {parent low};
    \draw[linkblue!80] (-1.60,-.62)--(-1,-.18);
  \end{scope}
\end{tikzpicture}
\caption{Schematic localization neighborhoods in one position coordinate.
At the marked child time bin, the intermediate neighborhood centered at the
selected low child \(b\) contains every low child of the parent in that bin
and is contained in the parent's high neighborhood. The angular bounds in
the actual neighborhoods are not depicted. These inclusions are used for
counting paths; the analytic evolution uses the exact coefficient maps with
vertex masks.}
\label{fig:adjacency-radii}
\end{figure}
